% !TEX program = xelatex
\documentclass[12pt,a4paper]{ctexart}
\usepackage[top=2.2cm, bottom=2.2cm, left=2.5cm, right=2.5cm]{geometry}
\usepackage{booktabs}
\usepackage{tabularx}
\usepackage{xcolor}
\usepackage{listings}
\usepackage{url}
\usepackage{fancyhdr}
\usepackage{lastpage}

\newcolumntype{L}{>{\raggedright\arraybackslash}X}

\lstset{
    basicstyle=\small\ttfamily,
    breaklines=true,
    frame=single,
    backgroundcolor=\color{gray!5},
    rulecolor=\color{gray!40},
    xleftmargin=0.8em,
    xrightmargin=0.8em,
    aboveskip=0.5em,
    belowskip=0.5em,
    showstringspaces=false,
}

\pagestyle{fancy}
\fancyhf{}
\fancyhead[L]{\small\textcolor{gray!50}{面向交通流的数据检测算法}}
\fancyhead[R]{\small\textcolor{gray!50}{环境说明}}
\fancyfoot[C]{\small\textcolor{gray!50}{\thepage\,/\,\pageref{LastPage}}}

\begin{document}

\begin{center}
    {\Large\bfseries 开发及运行环境说明}

    \vspace{0.4cm}
    面向交通流的数据检测算法 \quad·\quad 第十届全国高校智能交通创新创业大赛·企业赛道
\end{center}

\vspace{0.3cm}

\section{硬件环境}

\begin{table}[htbp]
    \centering
    \begin{tabularx}{\textwidth}{@{}p{3cm}L@{}}
        \toprule
        \textbf{项目} & \textbf{配置} \\
        \midrule
        开发机型 & Apple MacBook Pro（Apple M1） \\
        内存 & 16\,GB 统一内存 \\
        推理加速 & Apple MPS（Metal Performance Shaders），FP16 混合精度 \\
        最低运行要求 & $\geq$ 8\,GB 内存，Apple Silicon 或 NVIDIA GPU \\
        \bottomrule
    \end{tabularx}
\end{table}

\section{软件环境}

\begin{table}[htbp]
    \centering
    \begin{tabularx}{\textwidth}{@{}p{3cm}L@{}}
        \toprule
        \textbf{项目} & \textbf{版本 / 说明} \\
        \midrule
        操作系统 & macOS 26.5（Build 25F71），兼容 Windows 11 \\
        Python & 3.13.5（64-bit，ARM64 原生） \\
        IDE & JetBrains PyCharm Professional（教育版订阅） \\
        ultralytics & 8.4.51（YOLO 检测 + ByteTrack 跟踪） \\
        opencv-python & 4.x（视频读写、图像处理、可视化绘制） \\
        numpy & 2.x（矩阵运算、坐标转换） \\
        scipy & 1.x（空间距离、B 样条插值、JS 散度） \\
        hyperlpr3 & 最新版（车牌识别，可选依赖） \\
        openpyxl & 3.x（Excel 多 Sheet 报表导出） \\
        matplotlib & 3.x（轨迹图、验证报告图表） \\
        \bottomrule
    \end{tabularx}
\end{table}

推理设备切换：在 \texttt{src/config/settings.py} 中将 \texttt{DEVICE} 设为 \texttt{"mps"}（Apple Silicon）、\texttt{"cuda"}（NVIDIA GPU）或 \texttt{"cpu"} 即可。

\section{依赖安装}

\begin{lstlisting}[language=bash]
# 核心依赖（必需）
pip install ultralytics opencv-python numpy scipy openpyxl

# 可选依赖（推荐安装，提升数据完整性）
pip install hyperlpr3          # 车牌识别
pip install matplotlib         # 评测图表
pip install rapidocr-onnxruntime  # 辅助 OCR
\end{lstlisting}

\section{视频数据准备}

系统根据视频文件名自动识别进口方向，检测视频及模型权重须按以下规范放置。

\begin{table}[htbp]
    \centering
    \begin{tabularx}{\textwidth}{@{}p{3.2cm}L@{}}
        \toprule
        \textbf{启动方式} & \textbf{存放路径} \\
        \midrule
        源代码启动 &
        \texttt{src/assets/data/北进口\_20260420075959至20260420081500.mp4}\newline
        \texttt{src/assets/models/yolo26m.pt} \\
        \midrule
        编译程序 (exe) 启动 &
        \texttt{\_internal/src/assets/data/北进口\_20260420075959至20260420081500.mp4}\newline
        \texttt{\_internal/src/assets/models/yolo26m.pt} \\
        \bottomrule
    \end{tabularx}
\end{table}

视频命名格式：\texttt{\{进口方向\}\_\{年月日\}\{时分秒\}至\{年月日\}\{时分秒\}.mp4}，其中进口方向为“北进口”“东进口”“南进口”之一。三进口视频示例如下：

\begin{lstlisting}[language=bash]
北进口_20260420075959至20260420081500.mp4
东进口_20260420075958至20260420081459.mp4
南进口_20260420075959至20260420081500.mp4
\end{lstlisting}

\section{运行命令}

\begin{lstlisting}[language=bash]
# 交互式主入口
python3 -m src.main

# 核心流程：轨迹跟踪 + 车牌识别 + 断面过车 + 轨迹分组 + 报表
python3 -m src.main track

# 三进口批量处理
python3 -m src.main run-all

# 启动数据仪表盘（浏览器访问 localhost:8765）
python3 -m src.main dashboard
\end{lstlisting}

\subsection*{输出文件}

\begin{table}[htbp]
    \centering
    \begin{tabularx}{\textwidth}{@{}p{3cm}L@{}}
        \toprule
        \textbf{文件} & \textbf{内容} \\
        \midrule
        \path|outputs/trajectory.mp4| & 轨迹叠加可视化视频（1080p） \\
        \path|outputs/trajectory.csv| & 逐帧轨迹（坐标、速度、车道、车牌） \\
        \path|outputs/cross_section.csv| & 断面过车事件（流量、时距、速度、颜色） \\
        \path|outputs/trajectory_groups.csv| & 轨迹分组结果 \\
        \path|outputs/traffic_report.xlsx| & 多 Sheet 汇总报表 \\
        \bottomrule
    \end{tabularx}
\end{table}

\end{document}
